\chapter{Proof of Friedman's theorem}

\section{The free group and its reduced algebra}

\begin{definition}
  \label{def:free_group}
  We denote as $\free = \langle g_1, \ldots, g_r \rangle$ the \emph{free group} with $r$ generators.
\end{definition}

\section{Support of a real distribution}

\subsection{Stieltjes transform}

In this section we introduce a useful tool which will be essential to the next section.
Let $\RS$ denote the Riemann sphere.

\begin{lemma}
  \label{lem:Stieltjes_function}
  Let $z \in \RS$ such that $z \notin \R$.
  The function $x \in \R \mapsto \frac{1}{2i\pi} \frac{1}{x-z} \in \C$ is
  an element of $\smooth$.
\end{lemma}

\begin{proof}
  Usual functions.
\end{proof}

\begin{definition}
  \label{def:Stieltjes}
  \uses{lem:extension_csd,lem:Stieltjes_function}
  Let $\nu \in \csdistrib$.
  The \emph{Stieltjes transform} of $\nu$ is defined as the function
  $F_\nu : \RS - \R \rightarrow \C$ by $F(z) = \nu(x \mapsto \frac{1}{2 i \pi}\frac{1}{x-z})$.
\end{definition}

\begin{theorem}
  \label{thm:Stieltjes_holomorphic}
  \uses{def:Stieltjes}
  \notready
  Let $\nu \in \csdistrib$.
  The function $F_\nu : \RS - \R \rightarrow \C$ is analytic.
  It can be analytically extended to a $x \in \R$ if and only if $x \notin \supp \nu$.
\end{theorem}

\begin{proof}
  This might need some work!
\end{proof}

\subsection{Spectral radius and support}

We shall now show how one can access the support of a distribution
from its moments.

\begin{lemma}
  \label{lem:monomial_smooth}
  Let $p \geq 0$.
  The monomial function $x \mapsto x^p$ belongs in $\smooth$. Furthermore,
  for any integer $m \geq 0$, for any compact set $K$ of $\R$,
  if $K+ := \sup_{x \in K} |x|$, then
  \begin{equation*}
    \|x \mapsto x^p\|_{C^m(K)}
    \leq \max(p,1)^m \max (1, K_+)^p.
  \end{equation*}
\end{lemma}

\begin{proof}
  The regularity is well-known.
  Let $x \in K$ and $0 \leq k \leq m$ an integer.
  Then,
  the $k$-th derivative of the monomial at $x$ is given by
  $p(p-1) \ldots (p-k+1) x^{p-k}$ if $p\geq k$ and $0$ otherwise.
  In both cases,
  this is smaller in absolute value
  than $\max (p,1)^k \max(K_+,1)^p$ which is smaller than $\max(p,1)^m \max(K_+,1)^p$.
\end{proof}

\begin{definition}
  \label{def:moment}
  \uses{def:compactly_supported_distribution,lem:monomial_smooth}
  Let $\nu \in \csdistrib$.
  For $p \in \N$, we define the $p$-th moment $\nu_p := \nu(x \mapsto x^p)$.
\end{definition}

\begin{definition}
  \label{def:radius_csd}
  \uses{def:moment}
  Let $\nu \in \csdistrib$.
  We define $$\rho(\nu) := \limsup_{p \to \infty} |\nu_p|^{1/p} .$$
\end{definition}

\begin{lemma}
  \label{lem:radius_finite}
  \uses{def:radius_csd}
  Let $\nu \in \csdistrib$.
  Then $\rho(\nu) < \infty$.
\end{lemma}

\begin{proof}
  By Lemma \ref{lem:extension_csd}, there exists a compact set $K$,
  a constant $C$ and an integer $m$
  such that, for all $h \in \smooth$,
  $|\nu(h)| \leq C \|h\|_{C^m(K)}$.
Let $p \geq 0$ be an integer.
Then by Lemma \ref{lem:monomial_smooth}, if $K_+ := \sup_{x \in K} |x|$,
$$|\nu_p| \leq C \max(p,1)^m \max (K_+,1)^p$$
and hence $(|\nu_p|^{1/p})_{p \geq 0}$ is bounded.
\end{proof}

\begin{proposition}
  \label{prop:Stieltjes_expansion}
  \notready
  \uses{lem:radius_finite,def:Stieltjes}
  Let $\nu \in \csdistrib$.
  There exists $M > 0$ such that, for all $z \in \C$ satisfying $|z|>M$,
  \begin{equation*}
    F_\nu(z) = \sum_{p=0}^\infty \frac{\nu_p}{z^{p+1}}.
  \end{equation*}
\end{proposition}

\begin{proof}
  TODO
\end{proof}

The main result of this section is as follows.

\begin{theorem}
  \label{thm:support_limsup}
  \uses{lem:radius_finite,thm:Stieltjes_holomorphic}
  Let $\nu \in \csdistrib$.
  Then $\supp \nu \subseteq [-\rho(\nu), \rho(\nu)]$.
\end{theorem}

\begin{proof}
  Still need some work!
\end{proof}
